\documentclass[12pt]{article}

% ---- Packages ----
\usepackage[T1]{fontenc}
\usepackage{geometry}
\geometry{a4paper, margin=1.2in}
\usepackage{amsmath}
\usepackage{amssymb}
\usepackage{amsthm}
\usepackage{hyperref}
\usepackage{csquotes}
\usepackage[nameinlink,capitalise]{cleveref}
\usepackage{natbib}
\bibliographystyle{plainnat}
\setcitestyle{authoryear,open={(},close={)}}

% ---- Theorem environments ----
\newtheorem{theorem}{Theorem}[section]
\newtheorem{corollary}{Corollary}[theorem]
\newtheorem{proposition}{Proposition}[section]
\newtheorem{definition}{Definition}[section]
\newtheorem{remark}{Remark}[section]

% ---- Macros ----
\newcommand{\bind}{\phi}
\newcommand{\recomp}{\rho}
\newcommand{\Jsel}{J_{\mathrm{sel}}}
\newcommand{\Jclos}{J_{\mathrm{clos}}}
\newcommand{\Jmix}{J_{\mathrm{mix}}}
\newcommand{\Jgrowth}{J_{\mathrm{growth}}}
\newcommand{\dsem}{d_{\mathrm{sem}}}

% ---- Formatting ----
\setlength{\parskip}{0.5em}
\setlength{\parindent}{1.5em}
\sloppy

% ---- Title ----
\title{Constraint Regimes and the Structure of Surviving Objects\\
\large On the Invariance of Selection Pressure Across Censorship,
Feed Dynamics, and Cognitive Strategy}
\author{Flyxion}
\date{}

\begin{document}

\maketitle

\begin{abstract}
Three phenomena that appear historically and culturally unrelated---the
formal properties of Soviet-era science fiction, the structural dynamics
of feed-optimized media platforms, and the cognitive trajectory of agents
who voluntarily adopt high-discipline intellectual regimes---are shown to
be instances of a single underlying mechanism. In each case, a selection
environment imposes a fragment-level survivability condition, and the
global form of what survives is determined by that condition rather than
by the intentions of the agents involved. We formalize this through the
coupled variables of binding force $\bind(c)$---the degree to which a
piece of content retains semantic identity under arbitrary
decomposition---and recomposability $\recomp(c)$---the degree to which
its affective or functional charge survives context-switching. Systems
that select for high $\recomp$ converge to structurally shallow,
fragment-safe outputs. Systems that select for high $\bind$ converge to
internally coherent but low-circulation objects. Neither outcome depends
on stated intent. The invariant is the constraint structure. As a
corollary, voluntary adoption of a high-binding regime functions as an
exit from recomposability-dominated selection environments, and the
colloquial description of this strategy as ``galaxy brain'' discipline
masks a formal structural shift. We conclude that the full version of
this claim cannot itself propagate without transformation, which is not a
paradox but a fixed point of the framework.
\end{abstract}

\newpage

\tableofcontents

\newpage

% ============================================================
\section{Introduction: Three Phenomena, One Mechanism}
\label{sec:intro}
% ============================================================

Consider three situations that appear to have nothing in common.

First: a pair of writers in the Soviet Union produce science fiction
novels spanning decades. Their work is formally published, widely read,
and broadly regarded as subversive. It escapes censorship not by avoiding
critique but by being structured such that every extractable fragment can
be defended as allegory, genre fiction, or ambiguity. The corpus feels
``perfectly calibrated,'' as though a sophisticated editorial intelligence
had tuned each sentence to the threshold of permissibility.

Second: feed-optimized media platforms select for content that generates
response without requiring full processing---content that can be clipped,
quoted, embedded, and reacted to without loss of affective charge. The
discourse that survives this selection is increasingly characterized by
high affective intensity and low semantic stability: it circulates widely
but means something different in every context it enters. This is not
because the platform intended to destroy meaning. It is because the
platform optimized for a property that is approximately inverse to
meaning-preservation.

Third: an individual notices that their intellectual environment has
progressively supplied fewer genuinely difficult inputs---not because the
world has become less interesting, but because their social network can
no longer generate material above their internal processing threshold.
To continue developing, they shift from consuming externally curated
complexity to constructing it internally: writing formal papers, building
systems, enforcing compression, and selecting reading for maximal
transfer across domains. The output of this phase is structurally
different from the input phase: less recomposable, more coherent, lower
in circulation, higher in internal consistency.

These three situations look unrelated. One is about political censorship,
one about media economics, one about personal intellectual development.
But once they are expressed in terms of constraint structure rather than
domain content, they become instances of the same mechanism:

\begin{quote}
\emph{Any environment that imposes selection on fragments will shape the
internal structure of what survives. The trade-off between binding force
and recomposability determines whether surviving structures remain deep
or become shallow.}
\end{quote}

This paper develops that claim formally, traces it through each of the
three cases, and examines what follows from their unification.

% ============================================================
\section{The Core Variables}
\label{sec:variables}
% ============================================================

\begin{definition}[Binding Force]
\label{def:binding}
The binding force of a content object $c$, denoted $\bind(c) \in [0,1]$,
measures the degree to which the semantic content of $c$ is preserved
under arbitrary context transformation. Formally, let $T_\sigma$ denote a
context-transformation operator parameterized by context $\sigma$. Then:
\[
\bind(c) = 1 - \mathbb{E}_\sigma\bigl[\dsem(c,\, T_\sigma(c))\bigr]
\]
where $\dsem$ is a semantic distance measure and the expectation is taken
over the distribution of contexts into which $c$ may be reinserted.
High $\bind(c)$ indicates that meaning is preserved under
context-switching. Low $\bind(c)$ indicates that affective charge
survives while semantic content drifts.
\end{definition}

\begin{definition}[Recomposability]
\label{def:recomp}
The recomposability of $c$, denoted $\recomp(c) \in [0,1]$, measures the
degree to which the affective or functional charge of $c$ is preserved
under context transformation:
\[
\recomp(c) = \mathbb{E}_\sigma\bigl[A(T_\sigma(c))\bigr] \,/\, A(c)
\]
where $A(c)$ is the affective charge of $c$. Content is maximally
recomposable when it generates equivalent response regardless of the
context into which it is inserted.
\end{definition}

\begin{proposition}[Approximate Inversion]
\label{prop:inversion}
Under typical platform dynamics, $\recomp(c) \approx 1 - \bind(c)$.
More precisely:
\[
\recomp(c) \leq 1 - \bind(c) + \varepsilon(c)
\]
where $\varepsilon(c) \geq 0$ measures the degree to which the affective
charge of $c$ has become decoupled from its semantic content through
prior recomposition or deliberate design. Content with $\varepsilon(c)
= 0$ cannot be recomposed without losing both affect and meaning
simultaneously.
\end{proposition}

The quantity $\varepsilon(c)$ indexes what may be called
\emph{pseudo-binding}: the appearance of semantic constraint without
actual resistance to recomposition. Content that performs depth while
remaining affectively portable has high $\varepsilon(c)$. It is the
dominant rhetorical mode of feed-adapted discourse.

The inversion in \cref{prop:inversion} is the load-bearing element of
the entire framework. If increasing $\recomp$ necessarily decreases
$\bind$, then any selection environment that rewards $\recomp$ will
systematically erode $\bind$, regardless of what the agents within it
intend. No one has to decide to simplify things. The selection mechanism
does it automatically.

% ============================================================
\section{Selection Regimes and Closure Regimes}
\label{sec:regimes}
% ============================================================

\begin{definition}[Selection Regime]
\label{def:selreg}
A discourse or production environment operates in the \emph{selection
regime} when its dominant selection functional rewards recomposability:
\[
\Jsel(c) = \alpha \cdot E(c) + \beta \cdot \recomp(c) - \gamma \cdot K(c)
\]
where $E(c)$ is engagement, $\recomp(c)$ is recomposability, $K(c)$ is
cognitive cost to the consumer, and $\alpha, \beta, \gamma > 0$. Content
survives by being detachable from context. The stable equilibrium is
low-binding, high-recomposability discourse.
\end{definition}

\begin{definition}[Closure Regime]
\label{def:closreg}
An environment operates in the \emph{closure regime} when its dominant
selection functional rewards semantic invariance under transformation:
\[
\Jclos(c) = \lambda \cdot \bind(c) - \mu \cdot \delta(c)
\]
where $\bind(c)$ is binding force and $\delta(c)$ is detected
inconsistency under adversarial recomposition. Content survives by
remaining identical to itself across contexts. The stable equilibrium is
high-binding, low-entropy discourse.
\end{definition}

\begin{theorem}[Regime Incompatibility]
\label{thm:incompatibility}
For any content $c$:
\[
\nabla_c \Jsel \cdot \nabla_c \Jclos \leq 0
\]
with equality only on the set of measure zero where
$\bind(c) = \recomp(c) = \tfrac{1}{2}$. Gradient ascent under $\Jsel$
is gradient descent under $\Jclos$ and vice versa. There is no smooth
deformation that transforms a selection-regime system into a
closure-regime system.
\end{theorem}

\begin{proof}[Proof sketch]
$\Jsel$ increases in $\recomp(c) \approx 1 - \bind(c)$, so
$\nabla_c \Jsel$ points in the direction of increasing $\recomp$ and
decreasing $\bind$. $\Jclos$ increases in $\bind(c)$, so $\nabla_c
\Jclos$ points in the direction of increasing $\bind$ and decreasing
$\recomp$. These gradients are antiparallel.
\end{proof}

The two regimes are not points on a spectrum. They are characterized by
incompatible selection functionals, incompatible stable equilibria, and
incompatible populations of content. This is why attempts to operate
within a selection regime while preserving closure-regime
properties---the ``Trojan horse'' strategy---fail structurally rather
than contingently.

\begin{corollary}[Trojan Horse Bound]
\label{cor:trojan}
Let $c$ be content with binding force $\bind(c)$. Let $c'$ be its
translation into a more recomposable form with $\recomp(c') >
\recomp(c)$. Then:
\[
\bind(c') \leq \bind(c)
\]
with equality only if the translation is semantically lossless, which
requires $\recomp(c') = \recomp(c)$. Since the translation is performed
precisely to increase $\recomp(c')$, we have $\bind(c') < \bind(c)$
strictly. The propagating form binds less than the original.
\end{corollary}

The corollary removes a comforting belief: that ideas can be made widely
transmissible without alteration. Once the trade-off is formalized,
accessibility is not neutral. It is a transformation operator that
necessarily changes the object. The ``smuggled'' idea is not the same
idea; it is a projection that survives under higher recomposability
constraints.

% ============================================================
\section{Three Instances of the Same Structure}
\label{sec:instances}
% ============================================================

\subsection{Soviet-Era Fiction: Hostile Extraction as Selection Pressure}
\label{ssec:soviet}

The Strugatsky brothers published science fiction in the Soviet Union
across several decades. Their work was broadly regarded as subversive,
yet it escaped suppression. The standard explanation attributes this to
authorial craft: deliberate calibration of transgressive content to the
threshold of permissibility.

This explanation is correct in its description but wrong in its causal
structure. It implies a centralized optimizing agent---an author or a
regime monitoring discourse---producing the calibration through intent.
What actually produced it was a selection mechanism operating at the
fragment level.

The operative constraint was not: \emph{this book as a whole must be
defensible}. It was: \emph{any extractable fragment of this book,
removed from context, must be defensible as fiction, allegory, or
ambiguity}. This is a fragment-level survivability condition---exactly
the condition that maximizes $\bind(c)$ under hostile interpretation
rather than maximizes $\recomp(c)$ under feed amplification.

The result is a corpus where each piece is locally deniable under
arbitrary extraction. Every paragraph about a medieval tyrant can be
read as genre worldbuilding. Every passage about a population subjected
to daily seizures of manufactured euphoria can be read as speculative
fiction. The work feels ``perfectly calibrated'' not because it was
explicitly engineered that way, but because everything that was not
fragment-safe failed to survive.

\begin{remark}
The Soviet censorship regime and the feed selection regime impose
structurally identical fragment-level survivability conditions while
optimizing for opposite properties. Censorship selects for maximum
$\bind$ under hostile decomposition---every fragment must resist lethal
extraction. Feeds select for maximum $\recomp$ under algorithmic
amplification---every fragment must circulate without loss of charge.
Both produce corpora where no fragment is unsafe in the relevant sense.
Different penalty functions, same formal constraint structure, opposite
equilibria.
\end{remark}

The conspiracy version of this observation---that the authors were
managed by a state apparatus monitoring public discourse and serializing
feedback---is narratively satisfying but explanatorily redundant. The
selection mechanism produces the same output without coordination. This
is not to say the conspiracy version is false; it is to say it is a
limiting case of the selection mechanism, not the mechanism itself.

\subsection{Feed-Optimized Platforms: Recomposability as Fitness}
\label{ssec:feed}

Feed architectures optimize for a property of content that is not
identical to engagement but produces it: recomposability. A piece of
content is recomposable to the degree that it generates response without
requiring the respondent to have fully processed it. High recomposability
content can be clipped, quoted, embedded, and reacted to across contexts
without loss of affective charge.

The binding force of such content is correspondingly low. High-binding
content---content whose conclusion depends on the specific structure of
its premises, whose meaning changes when removed from context---is
maximally resistant to recomposition and therefore maximally penalized
by the selection functional $\Jsel$.

The practitioner class that has converged to high-$\recomp$ production
has not found a strategy for operating within a degraded environment.
It has converged to the only stable trajectory the environment admits.
This convergence does not require motivation or awareness. It requires
only that practitioners who do not converge lose visibility, and
practitioners who do converge retain it. Over time, the visible
population is the converged population.

\begin{theorem}[Convergence to Low-Binding Equilibrium]
\label{thm:convergence}
Under a platform selection functional that rewards recomposability, and
under boundary conditions that do not introduce strong countervailing
pressures, the visible discourse converges to a low-binding equilibrium
from which perturbations by high-binding content cannot produce
asymptotic stabilization.
\end{theorem}

\begin{proof}[Proof sketch]
By definition, $\Jsel$ assigns positive weight to $\recomp(c) \approx
1 - \bind(c)$. Content with high $\bind$ receives systematically lower
$\Jsel$ scores. Visibility is a monotone function of $\Jsel$. As
practitioners observe visibility distributions, adaptive pressure pushes
production toward low-binding strategies. A high-binding
injection---a single dense argument, a formal proof, a structurally
demanding essay---constitutes a local perturbation to the discourse
state, but does not alter the platform architecture or the boundary
conditions. The dynamics continue to be dominated by the gradients of
$\Jsel$, which re-amplify low-binding content and attenuate high-binding
content. The perturbation decays.
\end{proof}

\begin{corollary}[Insufficiency of Adaptive Strategies]
Adaptive strategies that reduce binding force to increase
recomposability do not alter the attractor structure of the system.
They accelerate convergence to the low-binding equilibrium.
\end{corollary}

This corollary reformulates the standard paradox of media critique: why
do sophisticated, well-intentioned practitioners keep producing content
that reinforces the systems they are analyzing? The structural
explanation is simpler than the standard invocation of co-optation or
false consciousness. The practitioners who remain visible are precisely
those whose output satisfied the recomposability constraint. The good
versions did not get suppressed in a conspiratorial sense. They failed
to meet the propagation threshold and became invisible. What looks like
ideological drift is survivorship bias over outputs.

\subsection{Voluntary Closure: The Self-Imposed Binding Regime}
\label{ssec:voluntary}

The third instance differs from the first two in that the selection
pressure is self-imposed rather than externally enforced. But the formal
structure is identical.

An individual operating within an externally curated intellectual
environment receives inputs filtered by other people's perception of
difficulty. This functions as a distributed selection mechanism: the
social network routes material above its own processing threshold toward
the individual perceived as capable of handling it. The incoming
material is already filtered for complexity. The individual's
intellectual development occurs within a propagation regime, except that
the selection criterion is perceived competence rather than
recomposability.

This arrangement saturates. Once the environment can no longer generate
inputs above the individual's internal binding threshold---once nothing
the network sends constitutes genuine difficulty---the distributed
intake system collapses. The environment routes fewer hard problems not
because the world has fewer hard problems, but because the social
network has exhausted its capacity to identify problems hard enough to
be useful.

Continuing development at this point requires switching regimes. The
individual must shift from consuming complexity to constructing it:
generating formal problems, enforcing compression across domains,
selecting reading for maximal transfer, producing objects with high
internal coherence. This shift increases $\bind$ at the expense of
$\recomp$. The outputs become less transmissible, more internally
structured, and more resistant to decomposition.

\begin{remark}
The saturation of the externally curated input channel is not a failure
state. It is a success metric: evidence that the first phase of
distributed intake operated correctly. The appropriate response is not
to seek out new sources of recomposable complexity but to switch
selection functionals---from $\Jsel$ to $\Jclos$, in effect. The second
phase is structurally harder because it removes the passive intake
channel. It is also the only phase that scales indefinitely, because
self-generated problems have no external ceiling.
\end{remark}

The individual in this third case is doing what the Strugatsky corpus
did unintentionally and what the feed resists structurally: constructing
objects whose internal coherence is not contingent on context. The
apparent ``galaxy brain'' quality---the tendency to see structural
patterns everywhere, to prefer domains where performance depends on
skill rather than chance, to read for transfer rather than
coverage---is the behavioral signature of operating under a
self-imposed closure regime. It is not a personality trait. It is a
constraint structure.

% ============================================================
\section{The Unified Principle}
\label{sec:unified}
% ============================================================

The three instances share a single formal property. In each case, a
selection environment imposes a condition at the fragment level, and
objects that satisfy the condition survive while objects that do not are
removed from the observable corpus. The global form of the surviving
corpus is determined by the fragment-level condition rather than by the
stated goals of the agents involved. The apparent calibration of the
output---its apparent fit to the selective environment---requires no
coordinating agent. It is a consequence of the filter.

Stated as a general principle:

\begin{quote}
\textbf{Constraint Invariance.} Any environment that imposes selection on
fragments will shape the internal structure of what survives. The
binding-recomposability trade-off determines whether surviving structures
remain deep or become shallow. The apparent intentions of the selecting
agent are irrelevant. The invariant is the constraint structure.
\end{quote}

This principle cuts through most of the usual explanatory apparatus in
media critique, political analysis, and accounts of intellectual
development. It does not require invoking coordination, conspiracy,
manipulation, or co-optation. It requires only specifying the selection
functional and reading off the equilibrium.

The unification also removes a structural asymmetry in how we typically
think about these cases. Soviet censorship is described as externally
imposed and coercive. Feed dynamics are described as emergent and
impersonal. Personal intellectual development is described as chosen and
autonomous. These descriptions are accurate about the surface mechanism.
They obscure the formal identity of the underlying structure. Specifying
the constraint reveals what the surface descriptions conceal.

% ============================================================
\section{Operationalization and Empirical Proxies}
\label{sec:measurement}
% ============================================================

The variables $\bind(c)$ and $\recomp(c)$ are defined formally in
\cref{def:binding,def:recomp}, but their usefulness depends on whether
they admit empirical approximation. While exact evaluation requires a
complete model of semantic transformation and affective response, both
quantities can be estimated through proxy procedures.

A fragment stability test samples contiguous fragments of $c$ of length
at least $k$ and evaluates interpretive variance across a distribution
of contexts $\sigma$. High-binding content exhibits low variance in
inferred meaning across contexts, while low-binding content exhibits
high variance. Letting $\mathcal{F}_k(c)$ denote the set of such
fragments, the expected variance over contexts provides a monotone proxy
for $1 - \bind(c)$.

Binding force can also be approximated through adversarial context
injection. Define an adversarial transformation $T_\sigma^{\mathrm{adv}}$
that places $c$ in a maximally hostile interpretive frame. The estimated
binding force is then:
\[
\widehat{\bind}(c) = 1 - \mathbb{E}_\sigma\bigl[\dsem(c,\,
T_\sigma^{\mathrm{adv}}(c))\bigr].
\]
Content that remains semantically stable under adversarial framing has
high estimated binding force.

Recomposability can be approximated by measuring the retention of
affective response under extraction. Let $A(f)$ denote measured
engagement or response to fragment $f$. Then:
\[
\widehat{\recomp}(c) =
\frac{\mathbb{E}_{f \in \mathcal{F}_k(c)}\bigl[A(f)\bigr]}{A(c)}.
\]
Content is highly recomposable when fragments produce comparable
responses to the original.

Finally, binding force may be estimated through semantic drift under
paraphrase. Let $P(c)$ denote a set of paraphrases of $c$. Then
$\widehat{\bind}(c) = 1 - \mathbb{E}_{p \in P(c)}[\dsem(c, p)]$.
High-binding content resists paraphrase without loss of meaning;
low-binding content is easily rephrased into semantically divergent
forms while retaining surface plausibility.

\begin{remark}
The existence of empirical proxies is sufficient to render the framework
falsifiable. If no observable procedure can distinguish between
high-binding and low-binding content, the variables are vacuous. If such
procedures exist---even approximately---the framework admits empirical
investigation without requiring full formal resolution of semantic
distance.
\end{remark}

% ============================================================
\section{Pseudo-Binding and the Simulation of Depth}
\label{sec:pseudobinding}
% ============================================================

The inversion relation in \cref{prop:inversion} admits an important
degeneracy: content can appear to exhibit high binding force while in
fact remaining highly recomposable. This phenomenon is captured by the
quantity $\varepsilon(c)$ and constitutes a distinct regime of
production.

\begin{definition}[Pseudo-Binding]
\label{def:pseudobinding}
A content object $c$ exhibits \emph{pseudo-binding} when it presents
surface features associated with high $\bind$---formal language,
hierarchical structure, technical vocabulary---while maintaining high
$\recomp$ under transformation. Formally, pseudo-binding is
characterized by $\recomp(c) \gg 1 - \bind(c)$ with $\varepsilon(c)$
large.
\end{definition}

Pseudo-binding arises when affective charge becomes decoupled from
semantic constraint. The content signals depth through stylistic markers
rather than structural necessity. Its apparent coherence does not depend
on the preservation of internal relationships between components, and
therefore survives extraction, paraphrase, and recomposition.

In selection-regime environments, content that mimics the surface
properties of high-binding objects while retaining recomposability gains
a selective advantage. It satisfies the engagement and portability
requirements of $\Jsel$ while borrowing credibility from the appearance
of structure. Over time, this produces a dominant rhetorical mode in
which depth is simulated rather than instantiated.

Pseudo-binding can be identified through adversarial decomposition. If
two fragments drawn from distinct regions of $c$ remain independently
interpretable and retain affective charge without requiring mutual
constraint, then the global structure is non-binding. The appearance of
coherence is not enforced by interdependence but assembled from locally
stable units. This distinguishes genuine high-binding content, where
removing a component degrades the remainder, from pseudo-bound content,
where components are substitutable.

\begin{proposition}[Dominance of Pseudo-Binding Under $\Jsel$]
\label{prop:pseudodominance}
In environments governed by $\Jsel$, pseudo-binding dominates over
genuine binding. For any content $c$ with high $\bind(c)$, there exists
a pseudo-bound approximation $c'$ such that $\recomp(c') > \recomp(c)$
and $\bind(c') < \bind(c)$, while preserving the appearance of depth.
Since $\Jsel$ rewards $\recomp$, $c'$ receives a higher selection score
and is preferentially amplified.
\end{proposition}

\begin{remark}
Pseudo-binding is the principal failure mode of attempts to operate
within selection regimes while preserving the appearance of closure. It
is not a transitional state but a stable attractor. Content that
successfully signals depth while remaining recomposable will outcompete
content that requires structural integrity to be understood.
\end{remark}

% ============================================================
\section{Mixed Regimes and Boundary Conditions}
\label{sec:mixed}
% ============================================================

The incompatibility result in \cref{thm:incompatibility} establishes
that pure selection and pure closure regimes are governed by opposing
gradients. However, most real-world environments are mixtures in which
multiple selection functionals operate simultaneously.

\begin{definition}[Mixed Regime]
\label{def:mixed}
A \emph{mixed regime} is an environment governed by a composite
selection functional:
\[
\Jmix(c) = \theta \Jsel(c) + (1-\theta)\Jclos(c),
\quad \theta \in [0,1].
\]
\end{definition}

For $\theta \approx 1$, the system is selection-dominated. For
$\theta \approx 0$, it is closure-dominated. The intermediate case
$\theta \approx \tfrac{1}{2}$ is unstable in practice, as small
perturbations in boundary conditions shift the effective gradient toward
one of the extremes.

In mixed environments, the observable corpus is determined not by the
full $\Jmix$ but by the component that controls visibility. Let $V(c)$
denote the visibility operator. If $V$ is monotone in $\Jsel$, then
even when $\Jclos$ is present, the visible discourse converges to the
selection-regime equilibrium. This explains why institutions that
nominally value coherence---academic journals, technical blogs, long-form
media---nevertheless exhibit selection-regime dynamics at the level of
visibility. High-binding work may exist within the system, but it does
not dominate the observable surface unless the visibility operator
itself is closure-aligned.

The effective regime is determined by boundary conditions rather than
formal objectives. When the cognitive cost of processing content is high
relative to available attention, $\Jsel$ dominates. When content is
subject to hostile extraction or sanction, $\Jclos$ dominates. When
distribution infrastructure rewards sharing and embedding, $\recomp$ is
amplified. When evaluation mechanisms assess content holistically---as
in formal proof verification or internal consistency checking---$\bind$
is amplified.

Mixed regimes can exhibit metastable configurations in which pockets of
high-binding content persist within a broader low-binding environment.
These pockets are typically shielded by access constraints, specialized
audiences, or alternative distribution mechanisms. They do not alter the
global attractor of the system.

\begin{proposition}[Visibility-Dominated Convergence]
\label{prop:visibility}
Let $V(c)$ be a visibility operator monotone in $\Jsel$. Then for any
mixed regime with $\theta > 0$, the visible corpus converges to the
selection-regime equilibrium, regardless of the magnitude of
$(1-\theta)\Jclos$.
\end{proposition}

\begin{remark}
Mixed regimes are often misinterpreted as evidence that the two regimes
can be reconciled. In practice, they demonstrate the opposite: the
dominant selection functional determines the observable equilibrium,
while the subordinate functional operates only within constrained
subspaces that do not scale to system-wide visibility.
\end{remark}

% ============================================================
\section{Failure Modes of High-Binding Regimes}
\label{sec:failure}
% ============================================================

The closure regime is not a neutral or universally superior alternative
to the selection regime. While it preserves semantic identity under
transformation, it introduces its own class of distortions. These
distortions arise from the same mechanism that produces its advantages:
the prioritization of internal coherence over external compatibility.

\begin{definition}[Over-Binding]
\label{def:overbinding}
A content system is \emph{over-bound} when increases in $\bind(c)$ no
longer correspond to increases in correspondence with external reality
or utility, but instead reinforce internal consistency alone.
\end{definition}

As $\bind(c) \to 1$, content becomes increasingly resistant to
transformation. This resistance can extend beyond protection against
semantic drift to resistance against legitimate revision. Highly bound
systems tend to reinforce their own internal constraints, leading to
closure in the sense of self-consistency rather than correctness. The
result is a system that cannot be wrong in the internal sense but may
be severely misaligned with external reality.

High-binding regimes can also reduce sensitivity to empirical signals.
Since evaluation occurs primarily through internal coherence,
discrepancies with external data may be reinterpreted rather than
incorporated. The tendency to identify structural invariants across
domains---which produces the transfer effects described in
\cref{ssec:voluntary}---can extend into overgeneralization when the
constraint to maintain coherence becomes stronger than the constraint to
discriminate between relevant and irrelevant patterns.

Operating under a closure regime also reduces exposure to recomposable
content streams. While this is a necessary condition for maintaining
high binding force, it reduces the diversity of incoming signals.
Without deliberate reintroduction of external constraints or adversarial
feedback, the system risks becoming locally coherent but globally
uninformative---an internally consistent model of a world it no longer
contacts.

\begin{proposition}[Trade-off Between Binding and Calibration]
\label{prop:calibration}
There exists a threshold $\tau$ such that for $\bind(c) > \tau$,
marginal increases in $\bind(c)$ produce diminishing or negative returns
in external calibration.
\end{proposition}

\begin{remark}
The failure modes of the closure regime mirror those of the selection
regime under inversion. Where selection regimes produce shallow but
adaptive systems, closure regimes produce deep but potentially rigid
ones. The symmetry is exact: each regime optimally avoids the failure
mode of the other while being maximally exposed to its own.
\end{remark}

% ============================================================
\section{Information-Theoretic Interpretation}
\label{sec:info}
% ============================================================

The variables $\bind(c)$ and $\recomp(c)$ admit an interpretation in
information-theoretic terms that clarifies the nature of the trade-off
and connects the formalism to existing theory.

Binding force can be interpreted as mutual information preserved under
transformation. Let $c$ be a content object with latent semantic
structure $S(c)$ and let $T_\sigma$ be a context transformation. Then:
\[
\bind(c) \;\approx\; \frac{I\bigl(S(c);\; S(T_\sigma(c))\bigr)}{H\bigl(S(c)\bigr)}.
\]
High-binding content preserves a large fraction of its semantic
information under transformation. Low-binding content exhibits rapid
information decay. Recomposability, by contrast, corresponds to the
preservation of low-order, high-amplitude signals---affective charge,
salience, recognizability---under projection, and is maximized when
these signals are largely independent of the higher-order structure of
$c$.

Context transformation $T_\sigma$ can be viewed as a projection operator
that maps $c$ into a lower-dimensional representation determined by the
target context. Under this view, recomposability corresponds to
robustness under lossy compression, while binding corresponds to
resistance to such compression. The inversion relation in
\cref{prop:inversion} is then a statement about compression trade-offs:
representations that preserve fine-grained structure are necessarily
more sensitive to projection, while representations that survive
projection must discard structure.

The pseudo-binding quantity $\varepsilon(c)$ corresponds in
information-theoretic terms to redundancy that preserves the appearance
of structure without increasing mutual information. In this sense,
pseudo-binding is high redundancy with low effective information content.
It survives projection because the redundant encoding carries the signal
even when structural constraints are removed.

\begin{proposition}[Rate-Distortion Interpretation]
\label{prop:ratedist}
Selection-regime environments minimize a distortion measure over
low-order signals subject to a constraint on representation cost, while
closure-regime environments minimize distortion over full semantic
structure without regard to representation cost.
\end{proposition}

\begin{remark}
The information-theoretic interpretation clarifies that the
binding-recomposability trade-off is not specific to discourse. It is a
general property of systems that must operate under transformation and
compression. The social, political, and cognitive examples developed in
this paper are instances of a broader class of information-processing
constraints.
\end{remark}

% ============================================================
\section{Visibility as a Selection Operator}
\label{sec:visibility}
% ============================================================

The analysis in previous sections implicitly assumes that survival is
equivalent to visibility. This section makes that assumption explicit by
introducing a visibility operator and clarifying its role in determining
the observable corpus.

\begin{definition}[Visibility Operator]
\label{def:visibility}
Let $V(c) \in [0,1]$ denote the visibility of content $c$. A visibility
operator is a mapping $V: \mathcal{C} \to [0,1]$ such that $V(c)$ is
monotone in the dominant selection functional of the environment. In
selection regimes, $V$ is monotone in $\Jsel$; in closure regimes, $V$
is monotone in $\Jclos$.
\end{definition}

Let $\mathcal{C}_{\mathrm{vis}} = \{ c \in \mathcal{C} \mid V(c) >
\theta_V \}$ denote the set of visible content above a threshold
$\theta_V$. The observable discourse is determined by
$\mathcal{C}_{\mathrm{vis}}$, not by the full population $\mathcal{C}$.

\begin{proposition}[Visibility-Induced Survivorship Bias]
\label{prop:survivorship}
If $V$ is monotone in $\Jsel$, then the observed distribution over
$\mathcal{C}_{\mathrm{vis}}$ is biased toward high-$\recomp$ and
low-$\bind$ objects, regardless of the distribution over $\mathcal{C}$.
\end{proposition}

This result provides a structural explanation for a common phenomenon:
the apparent drift of discourse toward simplified or distorted forms of
initially complex ideas. The drift is not necessarily due to changes in
belief or intent. It is a consequence of the visibility operator acting
on the population. What appears as corruption or co-optation of ideas
can often be reinterpreted as a visibility effect: high-binding variants
persist in $\mathcal{C}$ but fall below the visibility threshold, while
low-binding variants dominate $\mathcal{C}_{\mathrm{vis}}$. The
difference between existence and visibility is sufficient to produce the
observed pattern.

\begin{corollary}[Visibility as Effective Selection]
The composition $V \circ J$ defines the effective selection mechanism
of the system. Changing $J$ without changing $V$ does not alter the
observable equilibrium. Efforts to alter discourse by modifying content
without modifying the visibility operator are therefore structurally
limited.
\end{corollary}

% ============================================================
\section{Adversarial Optimization and Strategic Adaptation}
\label{sec:adversarial}
% ============================================================

The preceding analysis treats selection environments as passive filters.
In practice, many environments are adversarial: agents actively optimize
their outputs with respect to the selection functional. This section
extends the framework to such settings.

In selection regimes, adversarial optimization accelerates convergence
to high-$\recomp$ equilibria. Agents discover that reducing dependence
on context increases $\recomp(c)$, increasing affective intensity raises
$E(c)$, and lowering cognitive cost $K(c)$ increases accessibility.
Combined, these pressures produce content that is maximally extractable,
reactive, and context-invariant. Pseudo-binding emerges as a dominant
strategy under adversarial conditions.

\begin{proposition}[Adversarial Acceleration of Pseudo-Binding]
\label{prop:advpseudo}
In selection regimes, adversarial optimization over $\Jsel$ increases
$\varepsilon(c)$ over time. Agents converge to producing content that
maximizes apparent depth while preserving recomposability, because
genuine increases in $\bind(c)$ reduce $\recomp(c)$ and therefore
$\Jsel(c)$.
\end{proposition}

In closure regimes, adversarial optimization operates differently.
Agents seek to maximize $\bind(c)$ under adversarial testing, which
drives increased internal redundancy, formalization, and elimination of
interpretive ambiguity. This can produce over-binding as defined in
\cref{def:overbinding}: the regime's own adversarial pressure drives
$\bind(c) \to 1$, potentially beyond the calibration threshold $\tau$
identified in \cref{prop:calibration}.

Adversarial dynamics therefore amplify the characteristic failure modes
of each regime rather than introducing new equilibria. In selection
regimes, strategic optimization produces shallower and more
pseudo-binding content. In closure regimes, it produces more rigid and
less externally calibrated content. The direction of adversarial
pressure is set by the selection functional; the failure mode is a
consequence of following that pressure to its extreme.

\begin{remark}
This result clarifies that the stability of selection-regime equilibria
is not merely passive but is actively reinforced by the strategic
behavior of agents within the system. The convergence described in
\cref{thm:convergence} is accelerated, not merely sustained, by
adversarial adaptation.
\end{remark}

% ============================================================
\section{Classes of Transformation Operators}
\label{sec:transformations}
% ============================================================

The transformation operator $T_\sigma$ has thus far been treated
abstractly. A classification of common transformation types and their
effects on $\bind(c)$ and $\recomp(c)$ makes the framework more precise
and clarifies why different environments impose structurally different
constraints.

Clipping selects a subcomponent of $c$ and is the dominant
transformation in feed environments. It preserves $\recomp(c)$ when
fragments retain affective charge independently, but reduces $\bind(c)$
unless the content is locally self-contained---that is, unless each
fragment enforces its own semantic constraints without reference to
surrounding material.

Paraphrase replaces $c$ with a semantically similar representation.
Its effect depends on the redundancy of $c$: high-binding content
resists paraphrase without semantic loss, because its meaning is
enforced by structural necessity rather than surface phrasing.
Low-binding content admits paraphrase with minimal change in
interpretation, and in many cases paraphrase improves clarity without
any detectable loss.

Hostile framing embeds $c$ in a context designed to invert or distort
its meaning. This is the critical operator in censorship regimes, and
the one against which the Strugatsky corpus was optimized. High-binding
content maintains semantic identity under such transformation; low-binding
content is easily reinterpreted, because its meaning was never
sufficiently enforced by internal structure to resist reassignment.

Summarization compresses $c$ into a shorter representation. It increases
$\recomp$ of the result by reducing cognitive cost, but decreases
$\bind$ unless the compression is lossless---which is rare in practice,
since lossy compression is typically the mechanism that increases
recomposability. Summarization is thus the primary transformation by
which high-binding content is converted into recomposable form for
distribution.

Domain translation maps $c$ into a different conceptual framework. Its
effect on binding depends on whether the structural relationships within
$c$ have analogues in the target domain. High-binding content may retain
structure across domains if the relationships are sufficiently general;
low-binding content often becomes trivial or incoherent under domain
change, because its apparent structure was specific to the original
context.

\begin{proposition}[Transformation Sensitivity Ordering]
\label{prop:transformorder}
For typical content, the expected degradation of binding force under
each transformation class satisfies: clipping $\geq$ summarization
$\geq$ paraphrase. Hostile framing can exceed all others in adversarial
settings.
\end{proposition}

% ============================================================
\section{Scale Effects and Phase Transitions}
\label{sec:scale}
% ============================================================

The behavior of selection and closure regimes depends critically on
scale. This section characterizes how increasing system size induces
phase transitions between regimes.

In small communities or tightly coupled environments, context is shared,
cognitive cost constraints are relaxed, and evaluation can occur
holistically. Under these conditions, $\Jclos$ can dominate, and
high-binding content is sustainable. Technical seminars, research
collaborations, and specialized institutions are examples: the audience
is small enough that shared context makes binding force achievable
without the excessive cost it would impose at scale.

As system size increases, shared context diminishes, cognitive cost
constraints tighten, and evaluation becomes fragmentary. These changes
increase the weight of $\recomp(c)$ in effective selection, driving
convergence toward selection-regime equilibria. The probability that any
given piece of content is consumed out of context approaches 1 as the
audience grows, and once fragment-level evaluation dominates, $\Jsel$
becomes the effective selection functional regardless of nominal intent.

\begin{theorem}[Scale-Induced Regime Shift]
\label{thm:scale}
There exists a critical scale $S_c$ such that for system size $S > S_c$,
the effective selection functional becomes dominated by $\Jsel$,
regardless of initial conditions.
\end{theorem}

\begin{proof}[Proof sketch]
As $S$ increases, the probability that content is consumed out of
context approaches 1. Fragment-level evaluation therefore dominates. Since
fragment-level evaluation rewards recomposability, $\Jsel$ becomes the
effective selection functional at scale.
\end{proof}

Near $S_c$, systems may exhibit mixed-regime characteristics as described
in \cref{sec:mixed}. Beyond $S_c$, the transition is effectively
irreversible without structural changes to the visibility operator or
evaluation mechanism. This explains why small groups can sustain
high-binding discourse, why large platforms converge to
recomposability-dominated content, and why attempts to scale
closure-regime systems fail without altering evaluation structure. Scale
acts as a control parameter for regime selection. The drift toward
recomposability in large systems is not a failure of intent but a
consequence of scaling constraints.

\begin{remark}
This result provides a structural account of why high-binding discourse
has historically been associated with small, resource-constrained
institutions---monasteries, academies, research laboratories---rather
than with mass media or large platforms. The association is not
incidental. It is a consequence of the scale threshold below which
closure-regime dynamics can be maintained.
\end{remark}

% ============================================================
\section{Self-Generated Curriculum as Optimization}
\label{sec:curriculum}
% ============================================================

The transition described in \cref{ssec:voluntary} can be formalized as
an optimization problem over sequences of content objects. Rather than
selecting individual items for immediate consumption or production, the
agent selects a trajectory that maximizes long-term increase in binding
force.

\begin{definition}[Curriculum]
\label{def:curriculum}
A \emph{curriculum} is an ordered sequence of content objects
$\mathcal{K} = (c_1, c_2, \dots, c_T)$ together with a sequence of
transformations applied to them. The state of the agent at time $t$ is
represented by an internal model $M_t$, and the agent's effective
binding force $\bind(M_t)$ is defined analogously to \cref{def:binding},
measuring the stability of its internal representations under
transformation and recomposition.
\end{definition}

A growth functional over curricula is defined as:
\[
\Jgrowth(\mathcal{K}) =
\sum_{t=1}^{T} \Delta \bind_t - \eta \cdot C_t,
\]
where $\Delta \bind_t$ is the increase in effective binding force of the
agent's internal model induced by processing $c_t$, $C_t$ is the
cognitive cost incurred at step $t$, and $\eta > 0$ is a weighting
parameter. An optimal curriculum satisfies
$\mathcal{K}^* = \arg\max_{\mathcal{K}} \Jgrowth(\mathcal{K})$.

In practice, this requires selecting content that is neither trivial
(producing no increase in $\bind$) nor incomprehensible (producing high
cost with negligible increase in $\bind$), but lies in a regime of
productive difficulty. This is the formal correlate of the empirical
observation that learning is maximized when the difficulty of inputs
slightly exceeds current processing capacity.

In early phases, $\mathcal{K}$ is partially determined by external
selection mechanisms---social routing, institutional curricula, or feed
dynamics. These mechanisms saturate when they fail to produce inputs
with sufficient $\Delta \bind_t$. At this point, the agent must
approximate the argmax over $\mathcal{K}$ internally. A key operation in
this regime is compression: transforming $c_t$ into a more compact
representation without loss of semantic structure, increasing the density
of binding force per unit of cognitive cost and facilitating transfer
across domains.

\begin{proposition}[Scaling of Self-Generated Curricula]
\label{prop:scaling}
Externally generated curricula have an upper bound determined by the
complexity of the environment's selection mechanism. Self-generated
curricula do not have this bound, as the agent can generate arbitrarily
complex content subject only to its current model capacity.
\end{proposition}

\begin{remark}
The shift to self-generated curriculum is not optional for continued
growth beyond the saturation point. It is a structural requirement. The
``galaxy brain'' characterization corresponds to an agent that has
transitioned to optimizing $\Jgrowth$ directly, rather than inheriting
$\mathcal{K}$ from external selection regimes.
\end{remark}

% ============================================================
\section{Corollary: The Galaxy Brain as Regime Exit}
\label{sec:galaxybrain}
% ============================================================

The colloquial term ``galaxy brain'' originates in internet culture as a
satirical description of reasoning that escalates through increasingly
elaborate inferential steps to conclusions that appear absurd. The meme
encodes a double register: sincere admiration for the capacity to operate
at high levels of abstraction, and ironic awareness of the risk of
decoupling from empirical feedback.

This double register is structurally exact. It describes the phenomenology
of operating under a self-imposed closure regime from the perspective of
someone observing from within a selection regime. The outputs appear
``galaxy-brained''---elaborate, internally consistent, disconnected from
the propagation economy---precisely because they have been optimized for
$\bind$ rather than $\recomp$.

The irony in the meme performs the same function as the science-fiction
framing in the Strugatsky corpus: it provides protective ambiguity.
``I am describing a cognitive strategy'' and ``I am mocking a cognitive
failure'' are simultaneously available interpretations. The content
survives under hostile extraction---it can be read either way---because
neither reading requires retracting the other.

\begin{proposition}[Galaxy Brain as Closure Strategy]
\label{prop:galaxybrain}
Voluntary adoption of behaviors associated with ``galaxy brain''
discipline---preference for strategy over chance, reading for transfer,
compression of ideas into dense representations, avoidance of
recomposability-optimized content domains---functions as a self-imposed
closure regime. The colloquial framing masks a formal structural shift:
from optimizing for propagation to optimizing for coherence. The
cognitive gains are a byproduct of sustained exposure to high-binding
problem classes rather than a direct effect of claimed identity.
\end{proposition}

The practical implication is not to adopt a self-flattering identity. It
is to change the selection functional governing input and output.
Specifically: select inputs by binding force rather than by
recomposability, produce outputs that require context to be evaluated
rather than outputs that circulate without it, and tolerate the
reduction in visibility that results.

The resulting corpus---formal papers in obscure repositories, dense
monographs with limited circulation, systems that require significant
prior investment to understand---is not a failure mode. It is the
signature of correct operation under a self-imposed closure regime. It
is also, by the Trojan Horse Bound, the only kind of output whose
binding force is preserved rather than sacrificed in the act of
production.

% ============================================================
\section{The Fixed Point}
\label{sec:fixedpoint}
% ============================================================

This paper cannot propagate in its current form without transforming
into something with lower binding force. Any version that achieves
sufficient circulation to be widely recognized will have undergone the
transformation $\bind(c') < \bind(c)$ that the Trojan Horse Bound
predicts. The viral summary will be recomposable. The bound version
will be invisible.

This is not a paradox. It is a fixed point of the framework applied
recursively. The paper predicts its own fate: wide readership is mild
evidence against the integrity of the transmitted version. The correct
publication strategy follows directly: maximize $\bind$, deposit in a
context accessible to closure-regime readers, and do not attempt to
optimize for propagation. A version optimized for propagation would
refute itself by existing.

The framework does not conclude with a reform proposal. Reform proposals
that circulate within selection-regime environments have been selected
for recomposability and have therefore already had their binding force
reduced. The alternative to the selection regime already exists. It is
structured by a different selection functional, it produces a different
kind of object, and it is invisible from within the propagation economy
because visibility within the propagation economy requires the property
that the closure regime, by its structure, refuses to provide.

The full version of this claim cannot propagate without changing form.
That is not a counsel of despair. It is the most compressed statement
of the argument.

% ============================================================
\section{Conclusion}
\label{sec:conclusion}
% ============================================================

We began with three apparently unrelated phenomena and arrived at a
single mechanism. Soviet-era science fiction, feed-optimized media
dynamics, and voluntary high-discipline intellectual development are all
expressions of the same underlying structure: selection environments
that impose fragment-level survivability conditions shape the global
form of surviving objects through the binding-recomposability trade-off.

The conceptual contribution is the removal of motive from the
explanatory apparatus. Whether the selecting agent is a state censor,
an algorithmic recommendation system, or an individual imposing
discipline on their own cognitive inputs, the behavior of the resulting
corpus is determined by the constraint structure, not by intent. This
replacement of intentional explanation with structural explanation is
not cynicism. It is precision.

The practical contribution is the identification of regime exit as an
available option. The selection regime is not the only stable
environment for intellectual production. The closure regime exists,
operates under different principles, and produces objects with different
properties. The transition requires accepting reduced
recomposability---reduced transmissibility, reduced visibility, reduced
circulation---in exchange for increased binding force. This is not a
sacrifice. It is a change of optimization target.

The framework eating itself---this paper predicting its own limited
propagation---is the cleanest confirmation that the mechanism is
correctly identified. An argument that predicts its own fate under the
conditions it describes is not self-undermining. It is self-consistent.

\end{document}
